\section{Conclusion}

The research currently establishes a working, reproducible empirical foundation for studying regime-dependent stock-return predictability. The pipeline pulls real CRSP monthly data through WRDS, constructs predictors and next-month targets using explicit calendar timing, labels volatility regimes from an expanding historical VIX median, estimates Elastic Net and XGBoost under forward-looking restrictions, stores genuine out-of-sample predictions, and evaluates cross-sectional portfolio sorts with explicit safeguards against degenerate prediction cross-sections.

The preliminary real-data analysis covers 83 out-of-sample formation months from January 2019 through November 2025 using the partial size-and-momentum predictor set. The earlier AAPL validation shows that neither model outperforms the stock-specific expanding historical-mean benchmark, although XGBoost produces modestly lower forecast errors than Elastic Net. The subsequent cross-sectional analysis sharpens the interpretation. The tuned Elastic Net assigns one identical forecast to all eligible stocks in every OOS month and therefore contains no usable cross-sectional ranking signal in the current specification. Those month-model groups are excluded from decile formation rather than converted into arbitrary portfolios through identifier-based tie-breaking.

XGBoost produces non-degenerate prediction cross-sections in all 83 OOS months, but its mean monthly Spearman rank correlation with realized next-month returns is slightly negative. Equal-weighted D10-minus-D1 returns are economically close to zero overall. The value-weighted spread is positive overall and is concentrated in LOW-VIX months, while the HIGH-VIX value-weighted spread is essentially zero. Thus the current evidence does not support a simple claim that return predictability is stronger during high-volatility periods. Instead, it indicates that any preliminary economic-value signal is sensitive to portfolio weighting and appears concentrated among larger stocks during low-volatility months.

These findings remain descriptive rather than final empirical conclusions. The next stages add book-to-market through Compustat/CCM, expand the predictor set and production sample, finalize delisting-return treatment, implement turnover and the frozen 0/50-basis-point transaction-cost robustness specification, and add formal HAC/Newey--West inference and additional robustness analysis. The central research question therefore remains open, but the project has progressed from software and single-stock validation to a validated cross-sectional economic-value experiment with clearly identified model limitations.
